\section{Desenvolviendo la serpiente} \label{matProbSnakes}
Se ha perdido Dagoberto, la serpiente de tu amigo Rigoberto. Afortunadamente en su búsqueda por la jungla lograron encontrar $t$ serpientes, pero ahora caen en otro problema, ¿Cuál es el querido Dagoberto?\\
\begin{wraptable}{r}{0.215\linewidth}
	\begin{tabular}{|p{0.5em}|p{0.5em}|p{0.5em}|p{0.5em}|p{0.5em}|p{0.5em}|}
		\hline
		i&$\rightarrow$&$\rightarrow$&$\rightarrow$&$\rightarrow$&\urCornerArrow\\
		\hline
		$\Rsh$&$\rightarrow$&$\rightarrow$&$\rightarrow$&\urCornerArrow&$\downarrow$\\
		\hline
		$\uparrow$&$\Rsh$&$\rightarrow$&f&$\downarrow$&$\downarrow$\\
		\hline
		$\uparrow$&\llCornerArrow&$\leftarrow$&$\leftarrow$&\lrCornerArrow&$\downarrow$\\
		\hline
		\llCornerArrow&$\leftarrow$&$\leftarrow$&$\leftarrow$&$\leftarrow$&\lrCornerArrow\\
		\hline
	\end{tabular}
	\caption{Cómo recorrer cada serpiente}
	\label{recorrerSerpientes}
\end{wraptable}
Por suerte tu amigo recuerda que el artículo $7$ de la Ley $21.020$ obliga que todas las serpientes de Chile tengan escrito su nombre a lo largo de su cuerpo.
Un gran problema es que las $t$ serpientes que encontraron están enredadas como espiral, y a tu amigo le dan miedo las serpientes (a pesar de tener una como mascota).\\
¡Recae en ti la responsabilidad de desenredar todas las serpientes, para ayudar a descifrar lo que tienen escrito en su cuerpo y así traer a Dagoberto devuelta a su hogar!\\
Comenzando desde la esquina superior izquierda de la representación matricial, cada serpiente enredada se debe recorrer en sentido horario (ver Cuadro \ref{recorrerSerpientes}).\\

\textbf{INPUT}\\
La primera línea contiene un entero $t$, tal que $1 \leq t \leq 2 \cdot 10^5$, la cantidad de serpientes por desenredar.
Por cada serpiente se entrega una primera línea con dos enteros $m$ y $n$, tales que $1 \leq n, m \leq 20$, las dimensiones del espiral que forma la serpiente.\\
Luego se entregan $m$ líneas con $n$ caracteres, la representación matricial de la serpiente enredada.\\

\textbf{OUTPUT}\\
Para cada serpiente se debe imprimir la cadena de $n \cdot m$ caracteres que tiene escrita la serpiente en su cuerpo desenredado.\\

\textbf{EJEMPLOS}\\
\sampleBoxes
{2\\
	4 5\\
	Hola\_\\
	berts\\
	o!!oo\\
	gaD\_y\\
	3 3\\
	Fed\\
	o.e\\
	cir
}
{Hola\_soy\_Dagoberto!!\\
	Federico.
}

\sampleBoxes
{1\\
	10 2\\
	so\\
	ey\\
	t.\\
	nl\\
	ea\\
	i.\\
	pm\\
	re\\
	ej\\
	so\\
	.r
}
{soy.la.mejor.serpiente}

\pagebreak
